% *****************************************************************************
% Copyright 2025 ULFPP

% Author: A.S.Grm
% Description: This is an official LaTeX ULFPP thesis disposition template.
% Version: 2.0
% Date: 25/11/2025
%
% ----------------------------------------------------------------------
%                             !!! POMEMBNO !!!
%       
%   Predloga je narejena za prevajanje s poljubnim prevajalnikom:
%        - pdflatex, xelatex, luatex
%   kjer je potrebno uporabiti biber ukaz za izdelavo bibliografije !!!
% ----------------------------------------------------------------------
%
%                *** Izdelava PDF/A dokument ***
%
%
% pdfLatex generira navaden PDF, ki ga je potrebno pretvorit v PFD/A
%
% pretvorba v PDF/A online: 
%           https://www.pdfforge.org/online/en/pdf-to-pdfa
% 
% validacija PDF/A: 
%           https://www.pdfforge.org/online/en/validate-pdfa
%
% *****************************************************************************


% *****************************************************************************
%
%              !!! IZBIRA USTREZNEGA PROGRAMA in/ali SMERI !!!
%
%  Najprej je potrebno določit deklaracijo študija v 
%
%    \documentclass[<program študija>,<smer>]{fppdisp}
%
% - zamenjava za <program študija> z ustreznim programom
% - če je potrebna smer pa še vpis smeri: SMER={splošna,vojaška}
%
%  primer: \documentclass[pomnav,SMER=splošna]{fppdisp} 
%
%  - pomnav:   pomorstvo vsš, 1. stopnja, navtika 
%      potrebno dodati še SMER={splošna ali vojaška} 
%   
%  - pomstr:   pomorstvo vsš, 1. stopnja, pomorsko strojništvo
%      potrebno dodati še SMER={splošna ali vojaška}  
%   
%  - pom2:     pomorstvo, 2. stopnja
%      izbriši opcijo SMER v \setkeys{fppdisp}
%
%  - promtpl:  promet uni, 1. stopnja, TPL
%      izbriši opcijo SMER v \setkeys{fppdisp}
%
%  - prompttl: promet vsš, 1. stopnja, PTTL
%      potrebno dodati še SMER={splošna ali vojaška}
%
%  - prom2:    promet, 2. stopnja
%      izbriši opcijo SMER v \setkeys{fppdisp}
%
% *** Pomorstvo ***
\documentclass[pomnav,SMER=splošna]{fppthesis}
%\documentclass[pomstr,SMER=splošna]{fppthesis}
%\documentclass[pom2]{fppthesis}
% *** Promet ***
%\documentclass[prompttl,SMER=splošna]{fppthesis}
%\documentclass[promtpl]{fppthesis}
%\documentclass[prom2]{fppthesis}


% *****************************************************************************
%  Dodatni potrebni paketn in nove funkcionalnosti
% 
%  WARNING: minimalno, da v OverLeaf prevede v kratkem času. Če boste dodali
%          preveč paketov, OverLeaf ne bo pravočasno prevedel!      
% 
% *****************************************************************************
\usepackage{enumitem}   % okolje za naštevanje
\usepackage{siunitx}    % dodatno okolje za pravilno pisanje števil z enoto

% *** Caption to be bold face ***
\usepackage[labelfont=bf]{caption}

% *** Serif font za poglavja ***
\usepackage{titlesec}
\titleformat{\section}
{\normalfont\sffamily\Large\bfseries}
{\thesection}{1em}{}
\titleformat{\subsection}
{\normalfont\sffamily\large\bfseries}
{\thesubsection}{1em}{}


% *****************************************************************************
%  Literatura
% *****************************************************************************
\usepackage{csquotes}   % biber potrebuje
\usepackage[style=numeric,backend=biber,sorting=none]{biblatex}

% potrebno za deljenje dolgih URL naslovov
\setcounter{biburllcpenalty}{7000}
\setcounter{biburlucpenalty}{8000}

% ime datoteke, ki vsebuje navedeno literaturo
\addbibresource{literatura.bib}


% *****************************************************************************
%  Metapodatki - Potreben vnos pravilnih podatkov. Vsi podatki so NUJNI!
% *****************************************************************************

% Zamenjajte z vašimi podatki, vsi podatki so potrebni!
\setkeys{fppthesis}{
	NASLOVSLO={Preverjanje točnosti simulacij plovbe},
	NASLOVANG={Checking the accuracy of navigation simulations},
	AVTOR={Janez NOVAK},
	STEVILKA=123456,
	MENTOR={doc.dr.~Miha BONČA},
	LEKTOR={prof.~Alenka LEKTOR},
	MESEC=November,
	LETO=2025
}

%  *** Zahvala, povzetek in ključne besede ***
% - zahvala: napišite komu bi se radi zahvalaili
% - povzetekSLO: napiše se povzetek v slovenščini
% - povzetekANG: prevod povzetka v slovenščini
% - kljucnebesedeSLO: navedite ključne besede v slovenščini
% - kljucnebesedeANG: navedite ključne besede v angleščini

\zahvala{Zahvaljujem se \dots}

\povzetekSLO{V povzetku na kratko opišemo vsebinske rezultate dela. Sem ne sodi
	razlaga organizacije dela -- v katerem poglavju/razdelku je kaj, pač pa le opis
	vsebine.}

\povzetekANG{Prevod slovenskega povzetka v angleščino.}

% ključne besede ločite z vejico
\kljucnebesedeSLO{morska hidrodinamika, CFD, primerjava upora}

\kljucnebesedeANG{marin hydrodynamics, CDF, resistance comparison}


% *****************************************************************************
% ZAČETEK VSEBINE
% *****************************************************************************
\begin{document}
	
	
% *********************
% *** Novo poglavje ***
% *********************
\section{Uvod}
\label{sec:Uvod}

Uvodni del naloge ponavadi vsebuje:

\begin{itemize}[nosep]
	\item splošen pregled področja naloge,
	\item osnovno literaturo področja naloge,
	\item kaj bomo v nalogi naredili,
	\item kako je naloga sestavljena.
\end{itemize}

Namen uvodnega dela je opisati ozadje diplomskega dela, kaj so glavna izhodišča in hipoteza. Nadaljujete s splošnim opisom pristopa reševanja naloge, kjer se sklicujete na vire \cite{molland2017ship, lang}, ki opisujejo določene postopke, rezultate, algoritme in podobno, iz katerih se boste navdihovali oziroma nadaljevali vaše raziskovalno delo v nalogi.

Po uvodu pričnemo z opisom \textit{Materialov in metod}, ki jih v nalogi uporabimo. Nato nadaljujemo z naslednjim poglavjem \textit{Rezultati}, kjer opišemo rezultate, ki smo jih v okviru izdelave naloge pridobili. Nato zaključimo s poglavjem \textit{Diskusija}, kjer na kratko opišemo prednosti in slabosti našega diplomskega dela, recimo navzkrižno pokomentiramo rezultate, primerjamo uporabo različnih metod in podobno. Po želji dodamo lahko še poglavje \textit{Zaključek}, kjer opišemo zaključne misli na naše delo in kaj bi priporočali za naprej.


% *********************
% *** Novo poglavje ***
% *********************
\newpage
\section{Materiali in metode}
\label{sec:Materiali}

V tem poglavju natančno opišete uporabljeno opremo in postopke dela z opremo. Mogoče v začetnem delu opišete teoretične osnovne principe, ki jih boste uporabljali. Recimo pri določanju upora ladje s pomočjo meritev modela v bazenu, se uporabi Frudejev princip, ki je bil opisan v \cite{froude1888resistance}. To je bil primer enostavnega citiranja v NUMERIC stilu.


% *** Podpoglavje ***
\subsection{Opis trenja}

S podpoglavjem opišemo ločene segmente, kot je recimo pri uporu ladje delitev na upora zaradi trenja in preostali upor. 


% *********************
% *** Novo poglavje ***
% *********************
\newpage
\section{Rezultati}
\label{sec:Rezultati}

Rezultati vsebujejo vse kar ste izmerili izračunali in podobno, vaši lastni dosežki. Rezultate je treba pravilno predstaviti in natančno opisati. Iz osnovnih rezultatov lahko z dodatnim procesiranje pridobite nove, ki dodatno opišete, recimo v podpoglavju \textit{Obdelava podatkov}.

% *** Podpoglavje ***
\subsection{Obdelava podatkov}

Tukaj opišem način obdelave podatkov in prikažem dobljene rezultate ter jih natančno opišem.

\begin{table}[h]
	\caption{Primerjava med tovornima pristaniščema Koper in Reka.}
	\label{tab:primer_tabele1}
	\begin{center}
		\begin{tblr}{colspec={l|rr},
				%width=0\textwidth,
				row{even} = {white,font=\small},
				row{odd} = {bg=black!10,font=\small},
				row{1} = {bg=blue!20,font=\bfseries\small},
				hline{1} = {1pt,solid,black!60},
				hline{2} = {1pt,solid,black!60},
				hline{Z} = {1pt,solid,black!60},
				rowsep=3pt
			}
			% *** začetek podatko ***
			& Pretovor & TEU    & TEU/dvigalo & m3 \\ 
			Koper  & 353,880  &	88,470 & 594         & 12 \\
			Rijeka & 168,777  &	42,194 & 328         &	6
			% *** konec podatkov ***
		\end{tblr}
	\end{center}
	\begin{flushright}
		Vir:(Beškovnik in Twrdy, 2009)
	\end{flushright}
\end{table}


% *********************
% *** Novo poglavje ***
% *********************
\newpage
\section{Diskusija in zaključek}
\label{sec:Diskusija}

V tem poglavju opravi z diskusijo opravljenega dela. Recimo opišem, če so bili kje neznani problemi, zakaj so rezultati taki kot so, ali je mogoče kakšno vidno odstopanje in opišete razmišljanje zakaj je tako in podobne stvari.

V zaključku navedete na kratko, kaj je bilo narejeno in zaključite kaj bi lahko naredili bolje, kje vidite možnost nadaljevanja, uporabo drugih principov, kot so bili uporabljeni in podobno.



% ******************
% *** Literatura ***
% ******************
\cleardoublepage
\singlespacing
\addcontentsline{toc}{section}{Literatura}
\printbibliography

\end{document}
